\section{専門科目}
\subsection{問1}
まずスカラー行列$A = \lambda I$ ($\lambda\in \mb{F}_{p}^{\times})$の場合は
$A^{p-1} = I$を満たす. \aka{ $A\in \mr{SL}_{2}(\mb{F}_p)$なので$\lambda = \pm 1$. }$p\neq 2$だから, スカラー行列のときは2個ある. 以降$A$は非スカラー行列とする. 次を示す.
\lefba{
\begin{claim}
非スカラー行列$A\in \mr{SL}_{2}(\mb{F}_{p})$に対し,
\[A^{p-1} = I \iff Aの固有値は \lambda,\lambda^{-1}  かつ \lambda\in \mb{F}_{p}^{\times} \setminus{\pm 1}\]
\end{claim}
}
\prf
($\Rightarrow$) $A$の最小多項式$m(T)$は$T^{p-1}-1$を割り切る. $T^{p-1}-1$は重根を持たず, かつ$\mb{F}_{p}$で分解できる. $A$は非スカラー行列なので$m(T)$は1次ではない. $A$の固有値の一つを$\lambda\in \ovl{\mb{F}_{p}}^{\times}$とすれば,  $m(T)$は固有多項式$(T-\lambda)(T-\lambda^{-1})$に一致し, これが$T^{p-1}-1$を割り切るためには$\lambda\neq \lambda^{-1}$が必要だから$\lambda\neq \pm 1$.  $\lambda^{p-1}=1$なので$\lambda\in \mb{F}_{p}^{\times}$である. \\
($\Leftarrow$) 仮定より$A$は$P\in \mr{GL}_{2}(\mb{F}_{p})$を用いて対角化可能. $P^{-1}AP= \mr{diag}(\lambda, \lambda^{-1})$は明らかに$p-1$乗して$I$だから$A^{p-1}=I$. \qed \\
よってClaimの右の条件をみたす$A$の個数がわかればよい. 証明で$A$が$\mr{GL}_{2}(\mb{F}_{p})$の中で対角化可能であることが分かっているから, $A,B\in \mr{SL}_{2}(\mb{F}_{p})$が同じ固有値$\lambda,\lambda^{-1}\in \mb{F}_{p}^{\times} \setminus{\set{\pm 1}}$を持つならば, $A$と$B$は$\mr{GL}_{2}(\mb{F}_{p})$で共役である. したがって, 共役作用
\[\mr{GL}_{2}(\mb{F}_{p})\to \mr{Aut}(\mr{SL}_{2}(\mb{F}_{p})); \quad P\mapsto (A\mapsto P^{-1}AP)\]
において, $\sumib{\lambda & 0 \\ 0 & \lambda^{-1}}$の軌道
\[O_{\lambda} := \set{P^{-1}\sumib{\lambda & 0 \\ 0 & \lambda^{-1}} P \mid  P \in \mr{GL}_{2}(\mb{F}_{p}) } \subset \mr{SL}_{2}(\mb{F}_{p})\]
は固有値が$\lambda, \lambda^{-1}\in \mb{F}_{p}^{\times}\setminus{\set{\pm 1}}$であるような$A^{p-1} = I$を満たす$A\in \mr{SL}_{2}(\mb{F}_{p})$の全体の集合に等しい. このことと $O_{\lambda} = O_{\lambda^{-1}}$に注意すると, 非スカラー行列で$A^{p-1} = I$を満たすもの全体の個数は
\[\dfrac{1}{2} \sum_{\lambda \in \mb{F}_{p}^{\times}\setminus{\set{\pm 1}}}|O_{\lambda}| \]
に一致する. $\mr{GL}_{2}(\mb{F}_{p})$のうち$\sumib{\lambda & 0 \\ 0 & \lambda^{-1}}$を固定する固定部分群を$G_{\lambda}$とおくと, $|G_{\lambda}||O_{\lambda}| = |\mr{GL}_{2}(\mb{F}_{p})| = (p^2-p)(p^2-1)$が成り立つから,
上の式は
\[\dfrac{1}{2}\sum_{\lambda\in \mb{F}_{p}^{\times}\setminus{\pm 1} } (\mr{GL}_{2}(\mb{F}_{p}): G_{\lambda}) \]
である. そこで, $\lambda\in \mb{F}_{p}^{\times}\setminus{\set{\pm 1}}$に対する $|G_{\lambda}|$について調べよう. $P = \sumib{a&b\\ c&d}$が$P\in G_{\lambda}$を満たすことと
\[\bmat{a&b\\ c&d}\bmat{\lambda & 0 \\ 0 & \lambda^{-1}} = \bmat{\lambda & 0 \\ 0 & \lambda^{-1}}\bmat{a&b\\c&d} \]
を満たすことは同値で, これは
\[b\lambda^{-1} = b\lambda かつ c\lambda = c\lambda^{-1} \]
に同値である. $\lambda\neq \lambda^{-1}$であるため, $b=c=0$を得る. よって$|G_{\lambda}| = (p-1)^2$である. 以上より$G_{\lambda}$の指数は一定して$p(p+1)$だから, 非スカラー行列で$A^{p-1}=I$を満たすものの個数は
\[\dfrac{1}{2}p(p+1)|\mb{F}_{p}^{\times}\setminus{\set{\pm 1}}| = \dfrac{1}{2}p(p+1)(p-3)\]
以上より求める個数は
\[\bolm{\dfrac{1}{2}p(p+1)(p-3) + 2}\]

\subsection{問2}
(1): $e=u^2t^2-t$とおく. $f(X)$の根の一つは$\sqrt{-ut + \sqrt{e}}$である. $f(X)$は$\mb{C}[t]$における素元\footnote{$t$が素元であるのは$\mb{C}[t] \cong \mb{C}$が体で$(t)$が素イデアルだから.}$t$に関するアイゼンシュタイン多項式なので$\mb{C}[t] $上でも$\mb{C}(t)$上でも既約\footnote{Gaussの補題. UFDとその商体で成り立つ. }.  $M=\mb{C}(\sqrt{t},\sqrt{e})$とおく. 次を示す. \footnote{別証明: これにもKummer理論が使える. $R = \lan [t],[e]\ran \subset \mb{C}(t)^{\times}/(\mb{C}(t)^{\times})^2$で$[te]$が単位元でないので$[t]\neq [e]$. そして$[t],[e]\neq [1]$なので$R$の位数は4以上. よってKummer理論で$\mb{C}(t, \sqrt{R}).\mb{C}(t)$は4次以上. 4次以下は自明. }
\lefba{
\begin{claim}
$[M:\mb{C}(t)] = 4$.
\end{claim}
\prf $\mb{C}(\sqrt{t})/\mb{C}(t)$は明らかに2次. 明らかに$[M:\mb{C}(t)]\leq 4$だから, $<4$を仮定すると$M=\mb{C}(\sqrt{t})$である. $s:=\sqrt{t}$を変数と見る. $\sqrt{e}^2 - (u^2s^4 - s^2)$なので$\sqrt{e}$は$\mb{C}[s]$上整である. よって$\sqrt{e}\in \mb{C}[s]$だから$e=u^2s^4-s^2 = s^2(us-1)(us+1)$は$\mb{C}[s]$の平方元であるべきだがそれは$U\neq 0$ではありえない.\qed
}
$N=\mb{C}(\sqrt{e})$とおく. 明らかに$[N:\mb{C}(t)] = 2$である.
$N^{\times}/(N^{\times})^2$の部分群$\Delta = \lan [-ut + \sqrt{e}], [-ut - \sqrt{e}] \ran $を考える. もし$-ut + \sqrt{e}$が$(N^{\times})^2$に属すと, $\alpha = \sqrt{-ut + \sqrt{e}} \in N^{\times}$だがこれはおかしい($\mb{C}(\alpha)\subset N$だが左辺が$\mb{C}(t)$上4次). よって$\Delta$において二つの生成元は単位元ではない. \par
$N(\sqrt{\Delta}) = L = \mb{C}(\sqrt{ut-\sqrt{e}}, \sqrt{ut+\sqrt{e}})$である. Kummer理論により $\Delta^{\ast}:= \mr{Hom}_{Grp}(\Delta, \mb{C}_{1})$と$\gal{(N(\sqrt{\Delta})/N)}$は同型だから$|\Delta| = [L:N]$である. $f(X)$の既約性より$[-ut\pm \sqrt{e}]$は単位元ではない. また$[-ut + \sqrt{e}][-ut-\sqrt{e}] = [u^2t^2 - e] = [t]$が仮に単位元であるとすると$t\in (N^{\times})^2$だから$\sqrt{t} \in N$である. すると$N=\mb{C}(\sqrt{t},\sqrt{e}) = M$だがこれは拡大次数が異なりClaimに反する.\qed \\

(1)(別解)(Seeker): $f$の根の1つを$\alpha \in L$とする. $\mathbb{C}[t]$はUFDなので, 素元$t$についてのアイゼンシュタインの既約判定法より$f$は既約だとわかる. よって$[K(\alpha) : K]=4$. また, このとき$-\alpha \neq \alpha$で, $-\alpha$も$f$の根だし, $\pm \alpha$とは異なる根(標数0で$f$が既約なので重根はない)$\beta$をとると$-\beta$も$f$の根. よって$f = (X^2 - \alpha^2)(X^2 - t/ \alpha^2)$と因数分解できる. 係数比較して, $\alpha^2 + t/\alpha^2 = -2ut$. よって$t= \dfrac{\alpha^2}{2u+1/\alpha^2} = \dfrac{\alpha^4}{2u\alpha^2 + 1}$である. したがって$K(\alpha) = \mathbb{C}(t, \alpha) = \mathbb{C}(\alpha)$である. また, $L=K(\alpha , \beta)$となっている. あとは$\beta \notin K(\alpha)$を示せば$[L:K] = 8$だとわかる. そこで$\beta \in K(\alpha)=\mathbb{C}(\alpha)$と仮定して矛盾を導く. このとき$\beta = g(\alpha)/h(\alpha)$ ($g, h\in \mathbb{C}[\alpha]$)と表せる.  $\beta $は$X^2-t/\alpha^2$の根なので$\beta^2 -t/\alpha^2=0$であり,
    \[
    \frac{g^2}{h^2} = \frac{t}{\alpha^2}
    =\frac{\alpha^2}{2u\alpha^2 + 1}
    \]
となる. よって$g(\alpha)^2 (\sqrt{2u}\alpha + i)(\sqrt{2u}\alpha-i) = h(\alpha)^2\alpha^2$ ($\sqrt{2u}$は2乗して$2u$になる複素数の1つ). ここで$\mathbb{C}$が代数閉体で$\alpha \notin \mathbb{C}$だから$\alpha$は$\mathbb{C}$上超越的i.e.$\mathbb{C}[\alpha]$は$\mathbb{C}$上の1変数多項式環に同型. そこで先ほどの式の両辺を比較すると, 因数として$\sqrt{2u}\alpha +i$が左辺には奇数個含まれているのに, 右辺には偶数個含まれていることになり矛盾している. よって$\beta\notin K(\alpha)$であり, $[L:K] = 8$.

(2): ガロア群$G$は$f(X)$の根$\set{\alpha,-\alpha,\beta,-\beta}$の置換を定め, それで決まる. よって $G\subset S_4$とみなせる. $G$はSylow-2部分群で, Sylowの定理から非可換群$\lan (1234), (13)\ran$と(共役により)同型なので, $G$も非可換である.\qed


\subsection{問3}
(1): $1\in A'$と, $P,Q\in A'$に対して$P\pm Q \in A'$と$PQ\in A'$を示せばよい. $F = \mb{Z}_{\geq 0}^{n}$とし,  $I = (i_1,\dots, i_n) \in F$に対して, $x_I := x_1^{i_1}\dots x_{n}^{i_n}$とおく. $I\in F$に対し, $s(I) = \sum_{j} i_j$とする.
\[P = \sum_{I\in F} a_I x_I,\quad  Q=\sum_{I\in F} b_I x_I\]
とおく. $P=\phi(P), Q=\phi(Q)$とすると,
\[a_I = \zeta_m^{s(I)} a_I,\quad b_I = \zeta_m^{s(I)}b_I \]
である. $P+Q  = \sum_{I\in F} (a_I + b_I)x_I$だから,
\bealn{
\phi(P\pm Q) &= \sum_{I\in F} \zeta_m^{s(I)}(a_I + b_I) x_I \\
&= \sum_{I\in F} \zeta_m^{s(I)} a_Ix_I \pm \sum_{I\in F} \zeta_m^{s(I)}b_I x_I\\
&= \sum_{I\in F} a_Ix_I \pm \sum_{I\in F} b_Ix_I = P\pm Q
}
だから$P\pm Q \in A'$である. \par
$J,K \in F$に対して$J+K$を($\mb{Z}^n$の加法から)定める.
\[PQ = \sum_{I\in F} \parena{\sum_{J+K = I}(a_J b_K)}x_I\]
であり,
\bealn{
\phi(PQ) &= \sum_{I\in F} \zeta_m^{s(I)} \parena{\sum_{J+K=I} a_J b_K} x_I \\
&= \sum_{I\in F} \parena{\sum_{J+K=I} (\zeta_m^{s(J)}a_J)(\zeta_m^{s(K)}b_K) }x_I\\
&= \sum_{I\in F} \parena{\sum_{J+K=I} a_Jb_K} x_I \\
&= PQ
}
だから$PQ\in A'$. 最後に, $1\in A'$は$\phi(1) = 1$より明らか. \\
(2): $a_{p} = a_{q}=1$であるような$p\neq q$が存在したとする. このとき, $x_p^{m}\in A'$かつ$x_{p}^{m-1}x_q\in A'$である. $x_{p}^{m}$は$A'$で既約元であるが, 素元ではない. なぜなら, $x_{p}^{m-1}x_q \notin (x_{p}^{m})A'$だが, それの平方は$x_{p}^{2m-2} x_q^2 = x_{p}^{m}\cdot x_{p}^{m-2}x_{q}^2 \in (x_{p}^{m})A'$を満たすから. よって, $a_{p}=1$となる$p$は高々1個である. 逆にこのとき, もしすべての$a_{p}$が$0$なら$A'=A$だから, $A'$は多項式環に同型. そして, $a_{p} = 1$でそれ以外では$0$であるとき, $A' = \mb{C}[x_1,\dots, x_{p-1}, x_{p}^{m} , x_{p+1},\dots, x_{n}]$なので多項式環に同型である. 以上より,
\[\bolm{たかだか一つのpに対して a_p = 1}\]


\subsubsection{$a_i \leq 1$以外のとき}
$a_i\leq 1$以外のときに, 多項式環と同型になるパターンがある. たとえば, $m=6, n=2, a_1 = 3, a_2 = 2$とせよ. このとき$a_1i_1 + a_2i_2 \equiv 0\pmod{6}$であるような$i_1,i_2$の条件は$i_1\in 2\mb{Z}, i_2\in 3\mb{Z}$なので, $A' = \mb{C}[X^2, Y^3]$となる.

\subsection{問8}
$u_{\epsilon}, f \in C^2(\ovl{\Omega})$なので, $u_{\epsilon}, f \in L^1(\mb{R}^n)\cap L^2(\mb{R}^n)$である(ただし$u_{\epsilon}, f$は$\ovl{\Omega}$の外で0拡張する). 関数$f$のフーリエ変換
を次のように定義する:
\[\wiha{f}(\xi) :=\int_{\mb{R}^n} f(\xi)e^{-ix \cdot \xi} dx\quad (\xi\in\mb{R}^n)\]
ただし$x=(x_1,\dots ,x_n)\in \mb{R}^n$, $\xi = (\xi_1,\dots, \xi_n)\in \mb{R}^n$に対し$x\cdot \xi = \sum_{k=1}^{n} x_k\xi_k$である. $\Delta u_{\epsilon} = \sum_{k=1}^{n} \dfrac{\del^2}{\del x_k^2} u_{\epsilon}(x)$のフーリエ変換は部分積分より
{\small
\bealn{
\wiha{\Delta u_{\epsilon}}(\xi) &= \int_{[-R,R]^n} \parena{\sum_{k=1}^{n} \dfrac{\del^2}{\del x_k^2}u_{\epsilon}(x)} e^{-ix \cdot \xi} dx\quad (R>1)\\
&= \sum_{k=1}^{n} \parenb{\parenc{\dfrac{\del}{\del x_k}u_{\epsilon}(x)}_{x_k = -R}^{x_k = R} - \int_{[-R,R]^n}(-i\xi_k) \dfrac{\del}{\del x_k} e^{-ix\cdot \xi}dx } \\
&= \sum_{k=1}^{n} \parenb{(0-0) + i\xi_k \dfrac{\del}{\del x_k}e^{-ix\cdot \xi}dx} \\
&= \sum_{k=1}^{n} i\xi_k\int_{[-R,R]^n} \dfrac{\del}{\del x_k}u_{\epsilon}(x) e^{-ix\cdot \xi} dx\\
&= \sum_{k=1}^{n} (i\xi_k)^2 \int_{[-R,R]^n} u_{\epsilon}(x)e^{-ix\cdot \xi} dx\\
&= -|\xi|^2 \int_{\mb{R}^n}u_{\epsilon}(x) e^{-ix\cdot \xi} dx \\
&= -|\xi|^2 \wiha{u_{\epsilon}}(\xi)
}
方程式$-\epsilon \Delta u_{\epsilon} + u_{\epsilon} = f$ ($x\in \Omega$)の両辺をフーリエ変換すると,
\[\epsilon |\xi|^2 \wiha{u_{\epsilon}}(\xi) + u_{\epsilon}(\xi) = \wiha{f}(\xi)\]
これを$\wiha{u_{\epsilon}}$について解くと
\[\bolm{\wiha{u_{\epsilon}}(\xi) = \dfrac{\witi{f}(\xi)}{1+\epsilon |\xi|^2 } }\]
}
$\nabla u_{\epsilon} = \parena{\dfrac{\del u_{\epsilon}}{\del x_1}(x),\dots, \dfrac{\del u_{\epsilon}}{\del x_n}(x)}$で,
\[\abs{\nabla u_{\epsilon}(x)}^2 = \sum_{k=1}^{n} \abs{\dfrac{\del u_{\epsilon}}{\del x_k}(x)}^2 \]

\bealn{
\norm{\nabla u_{\epsilon}}^2_{L^2(\Omega)} &= \int_{\Omega} |\nabla u_{\epsilon}(x)|^2 dx \\
&= \sum_{k=1}^{n} \int_{\Omega} \abs{\dfrac{\del u_{\epsilon}}{\del x_k}(x)}^2 dx \\
&= \sum_{k=1}^{n} \norm{\dfrac{\del u_{\epsilon}}{\del x_k}}^2_{L^2(\Omega)}
}
ここで,
\bealn{
\wiha{\dfrac{\del u_{\epsilon}}{\del x_k}}(\xi) &= \int_{\Omega} \dfrac{\del u_{\epsilon}}{\del x_k}(x) e^{-ix\cdot \xi}dx \\
&= (i\xi_k) \int_{\Omega}u_{\epsilon}(x) e^{-ix\cdot \xi}dx \\
&= i\xi_k \wiha{u_{\epsilon}}(\xi)
}
なので, プランシュレルの定理より,
\bealn{
&\norm{\nabla u_{\epsilon}}^2_{L^2(\Omega)} \\
&= \sum_{k=1}^{n} \norm{\dfrac{\del u_{\epsilon}}{\del x_k}}^2_{L^2(\Omega)} \\
&= \dfrac{1}{(2\pi)^n} \sum_{k=1}^{n} \norm{\wiha{\dfrac{\del u_{\epsilon}}{\del x_k}}}^2_{L^2(\Omega)} \\
&= \dfrac{1}{(2\pi)^n} \int_{\mb{R}^n} \sum_{k=1}^{n} |i\xi_k\wiha{u_{\epsilon}}(\xi)|^2 d\xi\\
&= \dfrac{1}{(2\pi)^n}\int_{\mb{R}^n} |\xi|^2|\wiha{u_{\epsilon}}(\xi)|^2 d\xi \\
&= \dfrac{1}{(2\pi)^n }\int_{\mb{R}^n}|\xi|^2\abs{\dfrac{\wiha{f}(\xi)}{1+\epsilon|\xi|^2}}^2
d\xi \\
&= \dfrac{1}{(2\pi)^n} \int_{\mb{R}^n} \dfrac{|\xi|^2}{(1+\epsilon|\xi|^2)^2}|\wiha{f}(\xi)|^2 d\xi
}
よって
\bealn{
&\epsilon\norm{\nabla u_{\epsilon}}_{L^2(\Omega)}^2 \\
&=\dfrac{1}{(2\pi)^n} \int_{\mb{R}^n} \int_{\mb{R}^n} \dfrac{|\xi|^2}{(1+\epsilon|\xi|^2)^2} |\wiha{f}(\xi)|^2d\xi \\
&\leq \dfrac{1}{4(2\pi)^n} \int_{\mb{R}^n} |\wiha{f}(\xi)|^2 d\xi \\
&= \dfrac{1}{4(2\pi)^n}\norm{\wiha{f}}^2_{L^2(\mb{R}^n)} \\
&= \dfrac{1}{4}\norm{f}_{L^2(\Omega)}^2
}
だから
\[\bolm{\sqrt{\epsilon} \norm{\nabla u_{\epsilon}}_{L^2(\Omega)} \leq \dfrac{1}{2}\norm{f}_{L^2(\Omega)}}\]
また,
\bealn{
&\epsilon\norm{\Delta u_{\epsilon}}_{L^2(\Omega)} \\
&\leq \dfrac{\epsilon}{(2\pi)^{n/2}} \norm{\wiha{\Delta u_{\epsilon}}}_{L^2(\mb{R}^n)} \\
&= \dfrac{\epsilon}{(2\pi)^{n/2}} \norm{|\xi|^2 \wiha{u_{\epsilon}(\xi)}}_{L^2(\mb{R}^n)}\\
&= \dfrac{1}{(2\pi)^{n/2}} \norm{\dfrac{\epsilon |\xi|^2}{1+\epsilon |\xi|^2} \wiha{f}(\xi)}_{L^2(\mb{R}^n)}\\
&\leq \dfrac{1}{(2\pi )^{n/2}} \norm{\wiha{f}}_{L^2(\mb{R}^n)} \\
&= \norm{f}_{L^2(\Omega)}
}
よって
\[\bolm{\epsilon\norm{\Delta u_{\epsilon}}_{L^2(\Omega)} \leq \norm{f}_{L^2(\Omega)}}\]
であるから示せた.\qed \\
(ii) $u_{\epsilon} - f = \epsilon \Delta u_{\epsilon}$より,
\bealn{
&\norm{u_{\epsilon} - f}_{L^2(\Omega)} \\
&= \epsilon \norm{\Delta u_{\epsilon}}_{L^2(\Omega)} \\
&= \dfrac{\epsilon}{(2\pi)^{n/2}} \norm{\dfrac{|\xi|^2}{1+\epsilon |\xi|^2} \wiha{f}(\xi)}_{L^2(\mb{R}^n)} \\
&= \dfrac{\epsilon}{(2\pi)^{n/2}} \parena{\int_{\mb{R}^n} \abs{\dfrac{|\xi|^2}{1+\epsilon |\xi|^2 }}^2|\wiha{f}(\xi)|^2 d\xi }^{1/2} \\
&\leq \dfrac{\epsilon}{(2\pi)^{n/2}} \parena{\int_{\mb{R}^n}||\xi|^2 \wiha{f}(\xi)|}^{1/2} \\
&= \dfrac{\epsilon}{(2\pi)^{n/2}} \norm{|\xi|^2 \wiha{f}(\xi)}_{L^2(\mb{R}^n)} \\
&= \epsilon \norm{\Delta f}_{L^2(\Omega)}
}
仮定より$f$は2階までの各偏導関数が$\ovl{\Omega}$上の連続関数であるため, それらは$\Omega$上で有界連続である. よって$\Delta f$も$\Omega$上で有界連続であり, $\Omega$は有界領域だから,  $\norm{\Delta f}_{L^2(\Omega)}$が有限であることが分かる. したがって
\[\epsilon \norm{\Delta f}_{L^2(\Omega)} \to 0\]
が示された. \qed
